% paralleles.tex — f164r, vv. 123–164 : loci paralleli. Généré par .travail/f164r_lp_redac2/rediger.py
% à partir des collectes vérifiées (paralleles_123-143.json, paralleles_144-164.json) et des parallèles
% d'Austin–Olson 2004 retrouvés dans les sources ; règles : edition/CONVENTIONS.md § 4
\lp{123 σέβομαι Λατώ τ᾽ ἄνασσαν]}{Ar.~\textit{Thesm.}~120 Λατώ τε κρούματά τ᾽ Ἀσιάδος ; Ar.~\textit{Thesm.}~117 ἕπομαι κλῄζουσα σεμνὸν}
\lp{123 σέβομαι]}{Σ Ar.~\textit{Thesm.}~123 τὸ σέβομαι κοινόν ; Ar.~\textit{Nub.}~293 καὶ σέβομαί γ᾽ ὦ πολυτίμητοι ; Aesch.~\textit{Pers.}~694–695 σέβομαι μὲν προσιδέσθαι, / σέβομαι δ᾽ ἀντία λέξαι ; Aesch.~\textit{Supp.}~222–223 πάντων δ᾽ ἀνάκτων τῶνδε κοινοβωμίαν / σέβεσθ᾽ ; Pi.~\textit{P.}~6.23–25 μάλιστα μὲν Κρονίδαν, / βαρυόπαν στεροπᾶν κεραυνῶν τε πρύτανιν, / θεῶν σέβεσθαι ; Aesch.~\textit{Supp.}~922 τοὺς ἀμφὶ Νεῖλον δαίμονας σεβίζομαι ; Hsch.~α~340 ἀγάσματα· σεβάσματα}
\lp{123 Λατώ]}{Pi.~fr.~59~Boeckh (Σ Ar.~\textit{Eq.}~1263) τί κάλλιον ἀρχομένοισιν ἢ καταπαυομένοισιν, ἢ βαθύζωνόν τε Λατὼ καὶ θοᾶν ἵππων ἐλάτειραν ἀεῖσαι ; Scol. Attic. (Ath.~15.50,~694c–696a) ἐν Δήλῳ ποτ᾽ ἔτικτε τέκνα Λατώ, / Φοῖβον χρυσοκόμαν, ἄνακτ᾽ Ἀπόλλω, / ἐλαφηβόλον τ᾽ ἀγροτέραν / Ἄρτεμιν ; Carm. Delph.~3 μέλπετε δὲ Πύθιον ἄνακτα χρυσάορον / Φοῖβον, ὃν ἔτικτε Λατὼ πέτρᾳ Δηλίᾳ ; \textit{h.Ap.}~14 χαῖρε, μάκαιρ᾽ ὦ Λητοῖ, ἐπεὶ τέκες ἀγλαὰ τέκνα ; Ar.~\textit{Av.}~870 καὶ Λητοῖ Ὀρτυγομήτρᾳ ; Eur.~\textit{Hec.}~458–461 ἔνθα πρωτόγονός τε φοῖ-/νιξ δάφνα θ᾽ ἱεροὺς ἀνέ-/σχε πτόρθους Λατοῖ φίλᾳ ὠ-/δῖνος ἄγαλμα Δίας ; Eur.~\textit{Ion}~921–922 ἔνθα λοχεύματα σέμν᾽ ἐλοχεύσατο / Λατὼ Δίοισί σε καρποῖς ; Sapph. ap. Ath.~13.28 (571b–e) Λατὼ καὶ Νιόβα μάλα μὲν φίλαι ἦσαν ἑταῖραι}
\lp{123 Λατώ τ᾽ ἄνασσαν]}{Ar.~\textit{Thesm.}~970–971 καὶ τὴν τοξοφόρον / Ἄρτεμιν ἄνασσαν ἁγνήν ; Eur.~\textit{IT}~1230 ὦ Διὸς Λητοῦς τ᾽ ἄνασσα παρθέν᾽ ; Ar.~\textit{Ran.}~386–387 Δήμητερ ἁγνῶν ὀργίων / ἄνασσα συμπαραστάτει ; \textit{h.Cer.}~492 πότνια, ἀγλαόδωρ᾽, ὡρηφόρε, Δηοῖ ἄνασσα ; Hsch.~ω~267 Ὦπι ἄνασσα ; Eur.~\textit{El.}~678 καὶ Γαῖ᾽ ἄνασσα, χεῖρας ᾗ δίδωμ᾽ ἐμάς ; Aesch.~\textit{Eum.}~235 ἄνασσ᾽ Ἀθάνα, Λοξίου κελεύμασιν ; Aesch.~\textit{Eum.}~288 χώρας ἄνασσαν τῆσδ᾽ Ἀθηναίαν ἐμοὶ ; Soph.~\textit{Aj.}~774 ἄνασσα, τοῖς ἄλλοισιν Ἀργείων πέλας / ἵστω ; \textit{Od.}~3.380 ἀλλὰ ἄνασσ᾽ ἵληθι, δίδωθι δέ μοι κλέος ἐσθλόν ; Pi.~\textit{I.}~5.6 διὰ τεάν, ὤνασσα, τιμάν ; Pi.~fr.~205.1–2 S.-M. (Clem.Al.~\textit{Strom.}~6.10.83) Ἀρχὰ μεγάλας ἀρετᾶς, / ὤνασσ᾽ Ἀλάθεια ; \textit{h.Cer.}~75 Ῥείης ἠυκόμου θύγατερ, Δήμητερ ἄνασσα ; Hsch.~α~4594 ἄνασσα· δέσποινα ; Hsch.~ε~6695 Εὐάνασσα· ἡ Δημήτηρ ; Hsch.~ε~7126 Εὐρυάνασσα· ἡ Ἥβη ; Hsch.~θ~22 θαλάμων ἄνασσα· Ἀφροδίτη ; Phot.~α~1527 ἄνακτες καὶ ἄνασσαι· Ἀριστοτέλης ἐν τῇ Κυπρίων πολιτείᾳ· „καλοῦνται δὲ οἱ μὲν υἱοὶ καὶ ἀδελφοὶ τοῦ βασιλέως ἄνακτες, αἱ δὲ ἀδελφαὶ καὶ γυναῖκες ἄνασσαι“ ; Phot.~α~1667 ἄνασσα· βασίλισσα ; \textit{idem} Suid.~α~2060}
\lp{124 κίθαριν]}{Pi.~\textit{P.}~1.1–2 χρυσέα φόρμιγξ, Ἀπόλλωνος καὶ ἰοπλοκάμων / σύνδικον Μοισᾶν κτέανον· τᾶς ἀκούει μὲν βάσις, ἀγλαΐας ἀρχά ; Pi.~\textit{P.}~5.63–65 πόρεν τε κίθαριν, δίδωσί τε Μοῖσαν οἷς ἂν ἐθέλῃ ; Timoth.~\textit{Pers.}~229–233 (PMG~791) νῦν δὲ Τιμόθεος μέτροις ῥυθμοῖς τ᾽ ἑνδεκακρουμάτοις κίθαριν ἐξανατέλλει, θησαυρὸν πολύυμνον οἴξας Μουσᾶν θαλαμευτόν ; \textit{h.Ap.}~131 εἴη μοι κίθαρίς τε φίλη καὶ καμπύλα τόξα ; \textit{h.Ap.}~188 αὐτίκα δ᾽ ἀθανάτοισι μέλει κίθαρις καὶ ἀοιδή ; Hsch.~α~7665 Ἀσιάς· ἡ κιθάρα διὰ τὸ ἐν Ἀσίᾳ εὑρῆσθαι ; Ammon.~\textit{Diff.}~271 κίθαριν καὶ κιθάραν διαφέρειν φησὶν Ἀριστόξενος (fr. 102 Wehrli) ἐν τῷ Περὶ ὀργάνου. κίθαριν μὲν γὰρ εἶναι τὴν λύραν ; \textit{Il.}~3.54 οὐκ ἄν τοι χραίσμῃ κίθαρις τά τε δῶρ᾽ Ἀφροδίτης ; \textit{Od.}~1.159 τούτοισιν μὲν ταῦτα μέλει, κίθαρις καὶ ἀοιδή ; \textit{Od.}~8.248 αἰεὶ δ᾽ ἡμῖν δαίς τε φίλη κίθαρίς τε χοροί τε ; Thgn.~773,~778–779 Φοῖβε ἄναξ … / τερπομένῳ κιθάρῃ τ᾽ ἀμφ᾽ ἐρατῇ θαλίῃ / παιᾶσίν τε χορῶν ἰαχῇσί τε σὸν περὶ βωμόν ; Hsch.~θ~287 θεόπαιστος· ὑπὸ θεοῦ παιόμενος. ἢ κιθάρα ὑπὸ θεοῦ ⟨ἐν⟩διδομένη ; Hsch.~κ~2632 κίθαρις· κιθάρα, κιθάρισμα ; Hsch.~λ~1429 λύρα· κιθάρα ; Hsch.~χ~777 χρυσάορ· χρυσόηλ[ι]ος. χρυσοκίθαρις ; Suid.~α~4147 καὶ Ἀσιάδος κρούματα, τῆς κιθάρας. οὕτως Ἀριστοφάνης ; Zen.~6.14 Τί σοι ὁ Ἀπόλλων κεκιθάρικεν; τὸ κεκιθάρικεν, οἷον, ἐμαντεύσατο· ὡς φησὶν Αἰσχύλος ἐν Αἴαντι Λοκρῷ ; Epich.~fr.~68.1~K.-A. (Ath.~9.408d) κιθάραι, τρίποδες, ἅρματα, τράπεζαι χάλκιαι ; Epich.~fr.~108.1–2~K.-A. (Ath.~4.183c) Σεμέλα δὲ χορεύει / καὶ ὑπαυλεῖ σφιν † σοφὸς κιθάραι παριαμβίδας ; Pl. Com. fr.~10~K.-A. (Poll.~10.190) δότω τὴν κιθάραν τις ἔνδοθεν / καὶ τοὐπιπόρπαμα ; Aesch.~\textit{Supp.}~681 ἄχορον ἀκίθαριν ; \textit{Il.}~13.731 ἄλλῳ δ᾽ ὀρχηστύν, ἑτέρῳ κίθαριν καὶ ἀοιδήν ; \textit{Od.}~1.153 κῆρυξ δ᾽ ἐν χερσὶν κίθαριν περικαλλέα θῆκεν ; Alc.~fr.~41.15~V. ] κίθαρις δ[ ; Hdt.~1.24.5 λαβόντα τὴν κιθάρην, στάντα ἐν τοῖσι ἑδωλίοισι διεξελθεῖν νόμον τὸν ὄρθιον ; Pl.~\textit{R.}~3.399d λύρα δή σοι, ἦν δ᾽ ἐγώ, καὶ κιθάρα λείπεται κατὰ πόλιν χρήσιμα ; Σ Pi.~\textit{P.}~5.87 πόρεν τε κίθαριν: τὴν κιθαριστικὴν ἐμπειρίαν ; Hsch.~α~4076 Ἀμφιον⟨ί⟩η· κιθάρα ; Hsch.~σ~1371 σοφιστήν· πᾶσαν τέχνην σοφίαν ἔλεγον, καὶ σοφιστὰς τοὺς περὶ μουσικὴν διατρίβοντας καὶ τοὺς μετὰ κιθάρας ᾄδοντας ; Hsch.~φ~776 φόρμιξ· κιθάρα ἡ τοῖς ὤμοις φερομένη ; Phot.~α~2956 Ἀσιάς· ἡ κιθάρα, ἀπὸ Ἀσίου τινός … ὅτι ἐκλήθη Ἀσιὰς ἀπὸ τῶν χρωμένων Λεσβίων, οἵτινες οἰκοῦσι πρὸς τῇ Ἀσίᾳ}
\lp{124 κίθαρίν τε ματέρ᾽ ὕμνων]}{Eur.~\textit{El.}~716–717 λωτὸς δὲ φθόγγον κελάδει / κάλλιστον, Μουσᾶν θεράπων ; Pi.~\textit{O.}~2.1 ἀναξιφόρμιγγες ὕμνοι ; Carm. Delph.~3 χρυσέα ἁδύθρους κίθαρις ὕμνοισιν ἀναμέλπεται ; Σ Ar.~\textit{Thesm.}~120–121 ἐπειδὴ εἶπε κάτω « κιθάραν ματέρα ὕμνων » ; Σ Pi.~\textit{O.}~2.1a ἢ οἱ ἀνασσόμενοι ὑπὸ φορμίγγων ὕμνοι· πρότερον γὰρ ἐνδίδωσι τὸ μέλος ὁ κιθαριστής, ἔπειτα ἡ ᾠδὴ λέγεται ; Eur.~\textit{Ion}~881–884 ὦ τᾶς ἑπταφθόγγου μέλπων / κιθάρας ἐνοπάν, ἅτ᾽ ἀγραύλοις / κέρασιν ἐν ἀψύχοις ἀχεῖ / μουσᾶν ὕμνους εὐαχήτους ; Terp.~fr.~5~Bergk (Str.~13.2.4) σοὶ δ᾽ ἡμεῖς τετράγηρυν ἀποστρέψαντες ἀοιδὴν / ἑπτατόνῳ φόρμιγγι νέους κελαδήσομεν ὕμνους ; Σ Pi.~\textit{O.}~2.1b οἱ τῆς φόρμιγγος ἀνάσσοντες. πρῶτον γὰρ τὰ ᾄσματα συντίθενται, εἶθ᾽ οὕτω πρὸς αὐτὰ καὶ ἡ λύρα ἁρμόζεται}
\lp{124 ματέρ᾽ ὕμνων]}{Pi.~\textit{O.}~8.1 μᾶτερ ὦ χρυσοστεφάνων ἀέθλων, Οὐλυμπία ; Aesch.~\textit{Sept.}~224–225 πειθαρχία γάρ ἐστι τῆς εὐπραξίας / μήτηρ, γυνὴ σωτῆρος ; Soph.~\textit{Aj.}~172–174 ἦ ῥά σε Ταυροπόλα Διὸς Ἄρτεμις — / ὦ μεγάλα φάτις, ὦ / μᾶτερ αἰσχύνας ἐμᾶς — ; Pi.~\textit{N.}~3.1 ὦ πότνια Μοῖσα, μᾶτερ ἁμετέρα, λίσσομαι ; Pi.~\textit{N.}~5.6 οὔπω γένυσι φαίνων τέρειναν ματέρ᾽ οἰνάνθας ὀπώραν ; Scol. Attic. (Ath.~15.50,~694c) Πλούτου μητέρ᾽ Ὀλυμπίαν ἀείδω / Δήμητρα στεφανηφόροις ἐν ὥραις ; \textit{h.Merc.}~429–430 Μνημοσύνην μὲν πρῶτα θεῶν ἐγέραιρεν ἀοιδῇ, / μητέρα Μουσάων ; Σ Pi.~\textit{N.}~3.1 μητέρα ἑαυτοῦ εἶπε τὴν Μοῦσαν ὁ Πίνδαρος, ὡς ἂν ἐπιπνεόμενος ἐκ τῶν Μουσῶν ; Σ Pi.~\textit{O.}~8.1b μητέρα τῶν ἀγώνων τὴν Ὀλυμπίαν, οἷον μητρόπολιν ; \textit{Il.}~2.696 Δήμητρος τέμενος, Ἴτωνά τε μητέρα μήλων ; Hsch.~υ~188 ὑμνεῖ· εὐλογεῖ, ᾄδει}
\lp{125 ἄρσενι βοᾷ δοκίμῳ]}{\textit{Test.} Suid.~α~4016 Ἄρσενι βοᾷ δοκίμῳ: τῇ θαυμαστῇ φωνῇ καὶ δεδοκιμασμένῃ ; Σ Ar.~\textit{Thesm.}~125 εὐτόνῳ. προεῖπε γὰρ ὅτι ἐκ παρθένων ἐστὶν ὁ χορός. ἢ ἄρσενι βοᾷ δοκίμῳ, ἐπεὶ Ἀπόλλων ἐστὶν ὁ κιθαρίζων. τῆς οὖν κιθάρας ἄρσενι βοᾷ δοκίμου οὔσης φῶς ἔσσυτο}
\lp{125 ἄρσενι βοᾷ]}{Soph.~\textit{Phil.}~1455 καὶ κτύπος ἄρσην πόντου προβολῆς ; Soph.~fr.~523.2–3~R. (Polyxena) ἄρσενας χοὰς / Ἀχέροντος ὀξυπλῆγας ἠχούσας γόους ; Pers.~6.3–4 mire opifex numeris veterum primordia vocum / atque marem strepitum fidis intendisse Latinae ; Pl.~\textit{Lg.}~2.669c ὥστε ῥήματα ἀνδρῶν ποιήσασαι τὸ χρῶμα γυναικῶν καὶ μέλος ἀποδοῦναι ; Pl.~\textit{Lg.}~7.802e τὸ δὴ μεγαλοπρεπὲς οὖν καὶ τὸ πρὸς τὴν ἀνδρείαν ῥέπον ἀρρενωπὸν φατέον εἶναι, τὸ δὲ πρὸς τὸ κόσμιον καὶ σῶφρον μᾶλλον ἀποκλῖνον θηλυγενέστερον ; Quint.~\textit{Inst.}~1.10.31 quae nunc in scenis effeminata et impudicis modis fracta non ex parte minima, si quid in nobis virilis roboris manebat, excidit ; Nonn.~\textit{D.}~1.504–506 Πιερίδας μὴ κτεῖνε χορίτιδας, ὄφρα καὶ αὐταὶ / Φοίβου κῶμον ἄγοντος … / θῆλυ μέλος πλέξωσιν ὁμόθροον ἄρσενι μολπῇ ; Hdt.~1.17.1 ἐστρατεύετο δὲ ὑπὸ συρίγγων τε καὶ πηκτίδων καὶ αὐλοῦ γυναικηίου τε καὶ ἀνδρηίου ; Arist.~\textit{Pr.}~11.16 οἱ ἄγονοι, οἷον παῖδες, γυναῖκες, καὶ οἱ ἤδη γέροντες καὶ οἱ εὐνοῦχοι, ὀξὺ φθέγγονται, οἱ δὲ ἄνδρες βαρύ ; Quint.~\textit{Inst.}~1.8.2 sit autem in primis lectio virilis … non tamen in canticum dissoluta nec plasmate (ut nunc a plerisque fit) effeminata ; Hsch.~θ~489 θηλυγενής· ὁ ἐναντίος τῷ ἀρρενωπῷ ; Hsch.~α~7123 ἀρρενωπός· φοβερός}
\lp{125 ἄρσενι]}{Phot.~α~2871 ἄρρεν, οὐκ ἄρσεν … Ὅμηρος διὰ τοῦ σ … καὶ οἱ τραγικοί. πολλὰ δὲ καὶ ἄλλα τῶν Ἀττικῶν ὀνόματα φυλάττονται οἱ τραγικοὶ ὡς ἰσχνότερα τῆς τραγῳδίας καὶ πολιτικώτερα ; Parm.~fr.~12.4–6~D.-K. πάντα γὰρ ⟨ἣ⟩ στυγεροῖο τόκου καὶ μίξιος ἄρχει / πέμπουσ᾽ ἄρσενι θῆλυ μιγῆν τό τ᾽ ἐναντίον αὖτις / ἄρσεν θηλυτέρῳ ; Hsch.~α~7427 ἄρρενα· ἄρσενα}
\lp{125 βοᾷ]}{Pi.~\textit{P.}~10.38–39 παντᾷ δὲ χοροὶ παρθένων / λυρᾶν τε βοαὶ καναχαί τ᾽ αὐλῶν δονέονται ; Ar.~\textit{Thesm.}~102–103 λαμπάδα κοῦραι ξὺν ἐλευθέρᾳ / πατρίδι χορεύσασθε βοάν ; Eur.~fr.~370.8–9~Kannicht (\textit{Erechth.}) Λίβυος ἀχάεντος [λω]τοῦ, / κιθαρίδος βοαῖς ; Eur.~fr.~752g.8–12~Kannicht (\textit{Hyps.}) μέσῳ δὲ παρ᾽ ἱστῶι / Ἀσιὰς ἔλεγον ἰήιον / Θρῇσσ᾽ ἐβόα κίθαρις {Ὀρφέως} / μακροπόλων πιτύλων / ἐρέταισι κελεύσματα μελπομένα ; Pi.~\textit{N.}~5.38 ἔνθα μιν εὔφρονες ἶλαι σὺν καλάμοιο βοᾷ θεὸν δέκονται ; Pi.~\textit{O.}~3.8 φόρμιγγά τε ποικιλόγαρυν καὶ βοὰν αὐλῶν ἐπέων τε θέσιν ; Aesch.~\textit{Sept.}~394 ὅστις βοὴν σάλπιγγος ὁρμαίνει μένων ; \textit{Il.}~18.495 αὐλοὶ φόρμιγγές τε βοὴν ἔχον ; Phot.~β~185 βοήν· μάχην. καὶ μεγάλην φωνήν}
\lp{125 δοκίμῳ]}{Pi.~\textit{N.}~3.10–11 ἄρχε δ᾽ οὐρανοῦ πολυνεφέλα κρέοντι, θύγατερ, / δόκιμον ὕμνον ; Ar.~\textit{Pax}~1030–1031 τόν γε σοφῇ δόκιμον / φρενὶ πορίμῳ τε τόλμῃ ; Aesch.~\textit{Pers.}~546–547 κἀγὼ δὲ μόρον τῶν οἰχομένων / αἴρω δοκίμως πολυπενθῆ ; Eur.~\textit{Supp.}~277 ὦ δοκιμώτατος Ἑλλάδι ; Aesch.~\textit{Pers.}~87 δόκιμος δ᾽ οὔτις ὑποστὰς / μεγάλῳ ῥεύματι φωτῶν ; Hsch.~δ~2126 δόκιμον· χρήσιμον, τέλειον}
\lp{126–127 τᾷ φάος ἔσσυτο δαιμονίοις ὄμμασιν]}{Ar.~\textit{Av.}~217–222 ἵν᾽ ὁ χρυσοκόμας Φοῖβος ἀκούων / τοῖς σοῖς ἐλέγοις ἀντιψάλλων / ἐλεφαντόδετον φόρμιγγα θεῶν / ἵστησι χορούς· διὰ δ᾽ ἀθανάτων / στομάτων χωρεῖ ξύμφωνος ὁμοῦ / θεία μακάρων ὀλολυγή ; Ar.~\textit{Lys.}~1284 ὃς μετὰ μαινάσι Βάκχιος ὄμμασι δαίεται ; Aesch.~\textit{Pers.}~150–151 ἀλλ᾽ ἥδε θεῶν ἴσον ὀφθαλμοῖς / φάος ὁρμᾶται μήτηρ βασιλέως ; \textit{h.Ap.}~204–206 οἳ δ᾽ ἐπιτέρπονται θυμὸν μέγαν εἰσορόωντες / Λητώ τε χρυσοπλόκαμος καὶ μητίετα Ζεὺς / υἷα φίλον παίζοντα μετ᾽ ἀθανάτοισι θεοῖσι ; Σ Ar.~\textit{Thesm.}~126 τῆς κιθάρας φῶς ἔσσυτο τοῖς σοφοῖς· ἀπὸ γὰρ τῆς κιθάρας ἐδιδάχθησαν οἱ παλαιοὶ ᾄδοντες καὶ γράφοντες ποιήματα. ἢ τῇ Λητοῖ λαμπρότατα ὄμματα λέγει ; \textit{Il.}~1.474 μέλποντες ἑκάεργον· ὃ δὲ φρένα τέρπετ᾽ ἀκούων ; Σ Ar.~\textit{Thesm.}~126 (suite) πάλιν δὲ φῶς ἐγένετο αὐτοῖς ἀπὸ τῆς ἡμετέρας φωνῆς· ἄνευ γὰρ κιθάρας ἢ ἄνευ ἀνδρὸς τοῦ διδάσκοντος οὐ μαθήσεταί τις}
\lp{126 τᾷ φάος ἔσσυτο]}{Aesch.~\textit{Pers.}~395 σάλπιγξ δ᾽ ἀϋτῇ πάντ᾽ ἐκεῖν᾽ ἐπέφλεγεν ; Aesch.~\textit{Sept.}~103 κτύπον δέδορκα· πάταγος οὐχ ἑνὸς δορός ; Soph.~\textit{OT}~186 παιὰν δὲ λάμπει στονόεσσά τε γῆρυς ὅμαυλος ; Soph.~\textit{OT}~473–475 ἔλαμψε γὰρ τοῦ νιφόεντος ἀρτίως φανεῖσα ; Pi.~\textit{I.}~4.43 κεῖνον ἅψαι πυρσὸν ὕμνων ; Pi.~\textit{O.}~4.8–10 Οὐλυμπιονίκαν δέκευ / Χαρίτων ἕκατι τόνδε κῶμον, / χρονιώτατον φάος εὐρυσθενέων ἀρετᾶν ; \textit{h.Ap.}~201–203 αὐτὰρ ὁ Φοῖβος Ἀπόλλων ἐγκιθαρίζει / καλὰ καὶ ὕψι βιβάς· αἴγλη δέ μιν ἀμφιφαείνει ; Pi.~\textit{O.}~9.22 μαλεραῖς ἐπιφλέγων ἀοιδαῖς ; Σ Pi.~\textit{O.}~9.34d μαλεραῖς ἐπιφλέγων ἀοιδαῖς: λαμπραῖς, κατὰ μετωνυμίαν}
\lp{126 φάος]}{Ar.~\textit{Ach.}~1184–1185 ὦ κλεινὸν ὄμμα νῦν πανύστατόν σ᾽ ἰδὼν / λείπω φάος γε τοὐμόν, οὐκέτ᾽ εἰμ᾽ ἐγώ ; Hsch.~φ~1119 φῶς· τὸ ὁρώμενον τὸ παρὰ τοῦ θεοῦ δημιουργηθέν. καὶ τὸ πῦρ. καὶ ἡ χαρά ; Ar.~\textit{Av.}~1748 ὦ μέγα χρύσεον ἀστεροπῆς φάος ; Ar.~\textit{Ran.}~1528–1529 ἀπιόντι ποιητῇ / ἐς φάος ὀρνυμένῳ δότε δαίμονες οἱ κατὰ γαίας ; Ar.~\textit{Eq.}~973 ἥδιστον φάος ἡμέρας ; Aesch.~\textit{Ag.}~575 ὡς κομπάσαι τῷδ᾽ εἰκὸς ἡλίου φάει ; \textit{Il.}~18.11 χερσὶν ὕπο Τρώων λείψειν φάος ἠελίοιο ; Thgn.~569 λείψω δ᾽ ἐρατὸν φάος ἠελίοιο ; B.~3.79–80 ὅτι τ᾽ αὔριον ὄψεαι / μοῦνον ἁλίου φάος ; PMG~690.9 (Corinn.?) ] ἱαρὸν φάος σελάνας πασα[ ; Hsch.~φ~156 φάος· φῶς, ὄρθρου ἀνατολή ; Hsch.~φ~157 φάος ἠελίοιο· ἡ αὐγή}
\lp{126 ἔσσυτο]}{Ar.~fr.~718~K.-A. (=~698~Kock ; Heph.~8.8) τίς ὄρεα βαθύκομα τάδ᾽ ἐπέσυτο βροτῶν ; Eur.~\textit{HF}~919–920 λέγε, τίνα τρόπον ἔσυτο θεόθεν ἐπὶ / μέλαθρα κακὰ τάδε ; Eur.~\textit{Hel.}~1162 τείχεα δὲ φλογμὸς ὥστε Διός ἐπέσυτο φλόξ ; \textit{Il.}~14.518–519 ψυχὴ δὲ κατ᾽ οὐταμένην ὠτειλὴν / ἔσσυτ᾽ ἐπειγομένη ; \textit{Il.}~2.809 πᾶσαι δ᾽ ὠΐγνυντο πύλαι, ἐκ δ᾽ ἔσσυτο λαὸς ; Pratin.~\textit{PMG}~708.3–4 (Ath.~14.617c) ἐμὲ δεῖ κελαδεῖν, ἐμὲ δεῖ παταγεῖν / ἀν᾽ ὄρεα σύμενον μετὰ Ναϊάδων ; Pi.~\textit{P.}~4.135 ἐσσύμενοι δ᾽ εἴσω κατέσταν ; Hsch.~ε~6351 ἔσσυτο· ὥρμησεν. ἔθανεν}
\lp{126–127 δαιμονίοις ὄμμασιν]}{\textit{Il.}~1.603–604 οὐ μὲν φόρμιγγος περικαλλέος ἣν ἔχ᾽ Ἀπόλλων, / Μουσάων θ᾽ αἳ ἄειδον ἀμειβόμεναι ὀπὶ καλῇ ; Hes.~\textit{Th.}~36–41 Μουσάων ἀρχώμεθα, ταὶ Διὶ πατρὶ / ὑμνεῦσαι τέρπουσι μέγαν νόον ἐντὸς Ὀλύμπου … γελᾷ δέ τε δώματα πατρὸς / Ζηνὸς ἐριγδούποιο θεᾶν ὀπὶ λειριοέσσῃ ; Hes.~\textit{Th.}~68–70 αἳ τότ᾽ ἴσαν πρὸς Ὄλυμπον ἀγαλλόμεναι ὀπὶ καλῇ, / ἀμβροσίῃ μολπῇ· περὶ δ᾽ ἴαχε γαῖα μέλαινα / ὑμνεύσαις ; Pi.~\textit{P.}~1.12 κῆλα δὲ καὶ δαιμόνων θέλγει φρένας, ἀμφί τε Λατοίδα σοφίᾳ βαθυκόλπων τε Μοισᾶν ; Aesch.~\textit{PV}~902–903 μηδὲ κρεισσόνων θεῶν / ἔρως ἄφυκτον ὄμμα προσδράκοι με ; Hsch.~δ~72 δαιμόνιον· θεῖον. μόρσιμον ; Hsch.~δ~76 δαιμονίην· μακαρίαν, θειοτάτην}
\lp{126–127 φάος … ὄμμασιν]}{Ar.~\textit{Nub.}~285–286 ὄμμα γὰρ αἰθέρος ἀκάματον σελαγεῖται / μαρμαρέαις ἐν αὐγαῖς ; Soph.~\textit{OT}~80–81 ὦναξ Ἄπολλον, εἰ γὰρ ἐν τύχῃ γέ τῳ / σωτῆρι βαίη λαμπρὸς ὥσπερ ὄμματι ; Mosch.~\textit{Eur.}~86 ὄσσε δ᾽ ὑπογλαύσσεσκε καὶ ἵμερον ἀστράπτεσκεν}
\lp{127 ὄμμασιν]}{Ar.~\textit{Thesm.}~665 πανταχῇ δὲ ῥῖψον ὄμμα ; Ar.~\textit{Thesm.}~958 κυκλοῦσαν ὄμμα χρὴ χοροῦ κατάστασιν ; Ar.~\textit{Ran.}~1354 δάκρυα δάκρυά τ᾽ ἀπ᾽ ὀμμάτων ; Ar.~\textit{Eccl.}~1 ὦ λαμπρὸν ὄμμα τοῦ τροχηλάτου λύχνου}
\lp{127 ἁμετέρας τε δι᾽ αἰφνιδίου ὀπός]}{\textit{Test.} Suid.~δ~891 Δι᾽ αἰφνιδίου ὀπός: τουτέστι τῆς ἐνθουσιαστικῆς ; \textit{idem} Suid.~ο~471 (praem. Ὀπός· καὶ ; τουτέστιν ; om. τῆς)}
\lp{127 αἰφνιδίου]}{Aesch.~\textit{PV}~554–555 τὸ διαμφίδιον δέ μοι μέλος προσέπτα / τόδ᾽ ἐκεῖνό θ᾽, ὅ τ᾽ ἀμφὶ λουτρὰ ; Pl.~\textit{Ep.}~7.341c–d ἐκ πολλῆς συνουσίας γιγνομένης περὶ τὸ πρᾶγμα αὐτὸ καὶ τοῦ συζῆν ἐξαίφνης, οἷον ἀπὸ πυρὸς πηδήσαντος ἐξαφθὲν φῶς ; Pl.~\textit{Ion}~534b κοῦφον γὰρ χρῆμα ποιητής ἐστιν καὶ πτηνὸν καὶ ἱερόν, καὶ οὐ πρότερον οἷός τε ποιεῖν πρὶν ἂν ἔνθεός τε γένηται καὶ ἔκφρων ; Σ Ar.~\textit{Thesm.}~127 ἀντὶ τοῦ ἐνθουσιαστικῆς ; Hsch.~α~6867 ἄπτερος· αἰφνίδιος. παρὰ Ὁμήρῳ ὁ προσηνὴς ἢ ταχύς. Αἰσχύλος Ἀγαμέμνονι ; Th.~2.61.3 δουλοῖ γὰρ φρόνημα τὸ αἰφνίδιον καὶ ἀπρόσδοκητον ; Hsch.~δ~1175 διαμφίδιον μέλος· ἀλλοῖον, διαπαντὸς κεχωρισμένον. ἀμφὶς γὰρ χωρίς. Αἰσχύλος Προμηθεῖ δεσμώτῃ ; Call.~\textit{Cer.}~59–60 οἳ μὲν ἄρ᾽ ἡμιθνῆτες, ἐπεὶ τὰν πότνιαν εἶδον, / ἐξαπίνας ἀπόρουσαν ; Hsch.~α~1161 ἀδόκητον· αἰφνίδιον ; Hsch.~α~2186 αἰφνιδίως· ἐξάπινα ; Hsch.~α~6868 ἀπτέρωτα· ταχέα, αἰφνίδια ; Hsch.~ε~3051 ἔνθους· ἐνθουσιαστής, μαινόμενος, πνευματοφόρος ; Hsch.~ε~4075 ἐπαγωγή· αἰφνίδιος … Κασάνδρᾳ}
\lp{127 ὀπός]}{Hsch.~ο~2062 ὄψ· ὄψις. ὀφθαλμός. ἢ φωνή ; Ar.~\textit{Pax}~400 ἴθ᾽, ἀντιβολῶ σ᾽, ἐλέησον αὐτῶν τὴν ὄπα ; Alcm.~\textit{PMG}~3 fr.~1.1–4 Ὀλ]υμπιάδες περί με φρένας / ]ς ἀοιδας / ]ω δ᾽ ἀκούσαι / ]ας ὀπός ; Thgn.~531–532 ὁππότ᾽ ἀκούσω / αὐλῶν φθεγγομένων ἱμερόεσσαν ὄπα ; Ar.~\textit{Pax}~804–805 πικροτάτην ὄπα γηρύσαντος ἤκουσ᾽ ; Eur.~\textit{Ion}~1204 ἔσεισε κἀβάκχευσεν, ἐκ δ᾽ ἔκλαγξ᾽ ὄπα ; \textit{Il.}~16.76 οὐδέ πω Ἀτρεΐδεω ὀπὸς ἔκλυον αὐδήσαντος ; \textit{Od.}~11.421 οἰκτροτάτην δ᾽ ἤκουσα ὄπα Πριάμοιο θυγατρός ; B.~1.77 ]σαίνουσ᾽ ὀπί ; Hsch.~μ~886 μέροπες· ἄνθρωποι· διὰ τὸ μεμερισμένην ἔχειν τὴν ὄπα, ἤγουν τὴν φωνήν ; Hsch.~ο~538 ὀλίγῃ ὀπί· βραχείᾳ φωνῇ ; Hsch.~ο~965 ὀπά· θυρὶς ἡ ἐκ τοῦ κεράμου, καὶ φωνήν ; Hsch.~ο~966 ὄπα γηρύοντες· φωνὴν φθεγγόμενοι}
\lp{127 ὧν χάριν]}{Ar.~\textit{Thesm.}~586 πρὸς ποῖον ἔργον ἢ τίνος γνώμης χάριν}
\lp{128 ἄνακτ᾽ ἄγαλλε Φοῖβον]}{Hsch.~α~3113 ἀλλ᾽ ἄναξ· ἐξόδιον κιθαρῳ⟨δῶν⟩ τοῦτο, καθάπ⟨ερ ῥαψ⟩ῳδῶν καὶ τὸ «νῦν ⟨δὲ θεοί⟩» ; Hsch.~α~3944 ἀμφὶ ἄνακτα· ἀρχὴ νόμου κιθαρῳδικοῦ ; Eur.~\textit{Andr.}~1002–1003 πικρῶς δὲ πατρὸς φόνιον αἰτήσει δίκην / ἄνακτα Φοῖβον ; Soph.~\textit{OC}~1006–1007 εἴ τις γῆ θεοὺς ἐπίσταται / τιμαῖς σεβίζειν ; Soph.~\textit{OC}~1556–1557 εἰ θέμις ἐστί μοι τὰν ἀφανῆ θεὸν / καὶ σὲ λιταῖς σεβίζειν ; Scol. Attic. (Ath.~15.50,~694c–696a) Φοῖβον χρυσοκόμαν, ἄνακτ᾽ Ἀπόλλω ; Eur.~\textit{El.}~196–197 ἀλλ᾽ εὐχαῖσι θεοὺς σεβί-/ζουσ᾽ ἕξεις εὐαμερίαν ; Hsch.~α~2336 ἀκερσεκόμης· ὁ ἄκαρτος τὰς κόμας, ἐξ οὗ δηλοῦται τὸ ἀπενθές. διὸ καὶ Φοῖβος λέγεται ὁ Ἀπόλλων, καθαρὸς ὢν παντὸς πάθους ; Hsch.~φ~678 φοιβᾶναι· λαμπρῦναι. μαντεύσασθαι. κοσμῆσαι. καθᾶραι. ἁγνίσαι. καὶ φοιβάσαι ὁμοίως ; Hsch.~φ~683 φοῖβος· καθαρός, λαμπρός, ἁγνός, ἀμίαντος ; \textit{idem} Phot.~φ~237 (om. λαμπρός ; ἁγνός· ἢ μάντις· καὶ ἀμίαντος· καὶ ὁ Ἀπόλλων)}
\lp{128 ἄγαλλε]}{\textit{Test.} Phot.~α~87 ἀγάλλεσθαι· χαίρειν, μεγαλοφρονεῖν, καλλωπίζεσθαι … φασὶ δὲ καὶ ἀγάλλει τοὺς θεοὺς καὶ ἀγαλοῦμαι καὶ „ἄγαλλε“ καὶ ἀγάλλου καὶ ἀγάλλεται καὶ ἀγάλλῃ. καὶ πολλὴ παρὰ τοῖς παλαιοῖς ἡ χρῆσις ; Phot.~α~164 ἀγῆλαι· τιμῆσαι. τῶν πάνυ δὲ Ἀττικῶν ἐστιν ἡ λέξις. καὶ ἀγήλω καὶ ἀγαλοῦμεν ἐρεῖς καὶ „ἄγαλλε“ καὶ προσαγαλεῖ τὸν θεόν, ἀντὶ τοῦ εὔξεται καὶ τιμήσει … τὸ δὲ ἀγῆλαι κωμῳδικὸν καὶ ἐγγὺς γλώττης ; \textit{idem} Synag. Σᵇ~α~145 (τιμῆσαι θεόν, ἀγλαΐσαι ; ἔστι δὲ ἡ λέξις τῶν πάνυ Ἀττικῶν ; προσαγήλει ; om. τὸ δὲ … γλώττης) ; Ar.~\textit{Pax}~396–397 καί σε θυσίαισιν ἱεραῖσι προσόδοις τε μεγάλαισι / διὰ παντὸς ὦ δέσποτ᾽ ἀγαλοῦμεν ἡμεῖς ἀεί ; Hermipp.~fr.~8~K.-A. φέρε νῦν ἀγήλω τοὺς θεοὺς ἰοῦσ᾽ ἐγὼ / καὶ θυμιάσω τοῦ τέκνου σεσωσμένου ; Hes.~\textit{Th.}~68 ἀγαλλόμεναι ὀπὶ καλῇ ; Hsch.~α~481 ἀγῆλαι· σεμνῦναι, ἀναθεῖναι, κοσμῆσαι. καὶ εἰς ἄγασιν ἀγαγεῖν ; Hsch.~α~485 ἀγήλω· κοσμήσω ; Eup.~fr.~131~K.-A. ἀναθῶμεν νῦν χἠμεῖς τούτοις τασδὶ τὰς εἰρεσιώνας / καὶ προσαγήλωμεν ἐπελθόντες. χαίρετε πάντες. :: δεχόμεσθα ; Theopomp. Com. fr.~48~K.-A. καί σε τῇ νουμηνίᾳ / ἀγαλματίοις ἀγαλοῦμεν ἀεὶ καὶ δάφνῃ ; Pl.~\textit{Lg.}~11.931a τῶν δ᾽ εἰκόνας ἀγάλματα ἱδρυσάμενοι, οὓς ἡμῖν ἀγάλλουσι καίπερ ἀψύχους ; Σ Eur.~\textit{Med.}~1027 Φρύνιχος ἀγῆλαι ἀντὶ τοῦ εὔξασθαι. καὶ παρὰ Ἀριστοφάνει· διὰ παντὸς, ὦ δέσποιν᾽, ἀγαλοῦμεν ἡμεῖς ἀεί … δηλοῖ καὶ τὸ τιμῆσαι ; Hsch.~α~263 ἄγαλμα· πᾶν ἐφ᾽ ᾧ τις ἀγάλλεται. οὐχ ὡς ἡ συνήθεια τὸ ξόανον ; Hsch.~α~351 ἄγασαι· ἀγάλλου, ⟨ἱλάσκου⟩ ; Phot.~α~86 ἀγάλλει· ποιεῖ, σκευάζει, κοσμεῖ, τιμᾷ, προσεύχεται ; Phot.~α~163 ἀγῆλαι· τιμῆσαι θεόν, ἀγλαΐσαι. Εὔπολις Δήμοις … Ἀριστοφάνης Εἰρήνῃ … Ἕρμιππος Ἀρτοπώλισι … Θεόπομπος Πηνελόπῃ ; \textit{idem} Suid.~α~217 (θεὸν ἀγλαΐαις ; om. Εὔπολις … Πηνελόπῃ) ; Eur.~\textit{Med.}~1026–1027 πρὶν λέκτρα καὶ γυναῖκα καὶ γαμηλίους / εὐνὰς ἀγῆλαι λαμπάδας τ᾽ ἀνασχεθεῖν ; Eur.~\textit{Tro.}~151–152 ποδὸς ἀρχεχόρου πληγαῖς Φρυγίους / εὐκόμποις ἐξῆρχον θεούς ; Pi.~\textit{N.}~5.43 ἤτοι μεταΐξαντα καὶ νῦν τεὸς μάτρως ἀγάλλει / κείνου ὁμόσπορον ἔθνος, Πυθέα ; Hsch.~α~257 ἀγαλλητόν· ἐφ᾽ ᾧ ἄν τις ἀγασθείη ; Hsch.~α~273 ἀγαλλόμεναι· σεμνυνόμεναι}
\lp{128 ἄνακτ᾽]}{Thgn.~773 Φοῖβε ἄναξ, αὐτὸς μὲν ἐπύργωσας πόλιν ἄκρην ; Ar.~\textit{Vesp.}~143 ἄναξ Πόσειδον τί ποτ᾽ ἄρ᾽ ἡ κάπνη ψοφεῖ ; Ar.~\textit{Vesp.}~438 ὦ Κέκροψ ἥρως ἄναξ ; Ar.~\textit{Pax}~442 μηδέποτε παύσασθ᾽ αὐτὸν ὦ Διόνυσ᾽ ἄναξ ; Ar.~\textit{Plut.}~438 ἄναξ Ἄπολλον καὶ θεοί, ποῖ τις φύγῃ ; Eur.~\textit{Cyc.}~599 Ἥφαιστ᾽, ἄναξ Αἰτναῖε ; Eur.~\textit{Med.}~759 ἀλλά σ᾽ ὁ Μαίας πομπαῖος ἄναξ ; \textit{Il.}~15.253 τὸν δ᾽ αὖτε προσέειπεν ἄναξ ἑκάεργος Ἀπόλλων ; Pi.~\textit{P.}~8.67–68 ἄναξ, ἑκόντι δ᾽ εὔχομαι νόῳ / κατά τιν᾽ ἁρμονίαν βλέπειν ; \textit{h.Ap.}~140 αὐτὸς δ᾽, ἀργυρότοξε, ἄναξ ἑκατηβόλ᾽ Ἄπολλον ; Hsch.~α~4376 ἄνακτα· βασιλέα}
\lp{129 χαῖρ᾽ … παῖ Λατοῦς]}{\textit{h.Ap.}~545 καὶ σὺ μὲν οὕτω χαῖρε, Διὸς καὶ Λητοῦς υἱέ ; \textit{idem} \textit{h.Merc.}~579 (Μαιάδος pro Λητοῦς) ; \textit{h.Hom.}~1.20–21 καὶ σὺ μὲν οὕτω χαῖρε, Διώνυσ᾽ εἰραφιῶτα, / σὺν μητρὶ Σεμέλῃ, ἥν περ καλέουσι Θυώνην ; Scol. Attic. (Ath.~15.50,~694c) σέ τε, παῖ Διός, Φερσεφόνη· / χαίρετον, εὖ δὲ τάνδ᾽ ἀμφέπετον πόλιν}
\lp{129 χαῖρ᾽]}{Ar.~\textit{Thesm.}~111–112 χαῖρε καλλίστας ἀοιδᾶς / Φοῖβ᾽ ἐν εὐμούσοισι τιμαῖς ; Zen.~5.99 Σὺν δὲ θεοὶ μάκαρες: τοῦτο ἐπιλέγουσιν οἱ ῥαψῳδοί· ὡς καὶ οἱ κιθαρῳδοὶ, Ἀλλ᾽ ἄναξ μάλα χαῖρε ; Ar.~\textit{Thesm.}~972 χαῖρ᾽ ὦ ἑκάεργε ; Cratin.~fr.~237.2~K.-A. (Heph.~13.1) χαῖρε δή, Μοῦσα· χρονία μὲν ἥκεις ; Cratin.~fr.~359~K.-A. (Heph.~10.3) χαῖρ᾽ ὦ χρυσόκερως βαβάκτα κήλων, / Πάν, Πελασγικὸν Ἄργος ἐμβατεύων ; Call.~\textit{Cer.}~134 χαῖρε θεὰ καὶ τάνδε σάω πόλιν ἔν θ᾽ ὁμονοίαι ; PMG~934.19–20 (Paean Erythr.) χαῖρέ μοι, ἵλαος δ᾽ ἐπινίσεο / τὰν ἀμὰν πόλιν εὐρύχορον ; PMG~935.25–26 χαῖρ᾽ ὦ μεγάλα ἄνασ-/σα Μᾶτερ Ὀλύμπου ; PMG~937.11–15 χαίρετε ἀθάνατοι πάντες θεοὶ αἰὲν ἐόντες / ἀθάναταί τε θεαὶ … / εὐμενεῖ σὺν ὄλβωι ; \textit{h.Ven.}~292 χαῖρε, θεά, Κύπροιο ἐυκτιμένης μεδέουσα ; Call.~fr.~112.8~Pf. χαῖρε, Ζεῦ, μέγα καὶ σύ, σάω δ᾽ [ὅλο]ν οἶκον ἀνάκτων}
\lp{129 ὄλβιε]}{Eur.~\textit{Hipp.}~1440 χαίρουσα καὶ σὺ στεῖχε, παρθέν᾽ ὀλβία ; Aesch.~\textit{Supp.}~524–526 ἄναξ ἀνάκτων, μακάρων / μακάρτατε καὶ τελέων / τελειότατον κράτος, ὄλβιε Ζεῦ ; Ar.~\textit{Lys.}~1286 πότνιαν ἄλοχον ὀλβίαν ; \textit{Od.}~24.36 ὄλβιε Πηλέος υἱέ, θεοῖς ἐπιείκελ᾽ Ἀχιλλεῦ ; Hsch.~ο~511 ὀλβιόδαιμον· τρισόλβιε. μακάριε, κατὰ τὴν εὐδαιμονίαν ; Sapph. ap. Heph. Ench.~15.26 ὄλβιε γαμβρέ, σοὶ μὲν δὴ γάμος, ὡς ἄραο ; Thgn.~165–166 οὐδεὶς ἀνθρώπων οὔτ᾽ ὄλβιος οὔτε πενιχρὸς / οὔτε κακὸς νόσφιν δαίμονος οὔτ᾽ ἀγαθός ; Hsch.~ο~512 ὄλβιος· πλούσιος, εὐδαίμων. ὁ διὰ τοῦ ὅλου βίου μακαριστός ; \textit{idem} Suid.~ο~134 (εὐδαίμων, μακάριος ; om. ὁ διὰ … μακαριστός) ; Hsch.~π~367 πανόλβιος· παμμακάριστος. ⟨ἢ ὁ τὰ κακὰ ἀποστρέφων⟩ ; Hsch.~τ~1429 τρισόλβιος· τρισμακάριστος ; Ammon.~\textit{Diff.}~351 ὄλβιος καὶ εὐτυχὴς διαφέρει. ὄλβιος μὲν γάρ ἐστιν ὁ τελείαν τὴν εὐδαιμονίαν ἔχων, εὐτυχὴς δὲ ὁ ζῶν ἡδέως καὶ ἀλύπως}
\lp{129 παῖ Λατοῦς]}{Eur.~\textit{IT}~1234 εὔπαις ὁ Λατοῦς γόνος ; Eur.~\textit{Ion}~125–127 ὦ Παιὰν ὦ Παιάν, / εὐαίων εὐαίων / εἴης, ὦ Λατοῦς παῖ ; Eur.~\textit{HF}~687–690 παιᾶνα μὲν Δηλιάδες / ὑμνοῦσ᾽ ἀμφὶ πύλας τὸν / Λατοῦς εὔπαιδα γόνον / εἱλίσσουσαι καλλίχορον ; Eur.~\textit{IT}~126–127 ὦ παῖ τᾶς Λατοῦς, / Δίκτυνν᾽ οὐρεία ; Eur.~\textit{Ion}~885 σοὶ μομφάν, ὦ Λατοῦς παῖ}
\lp{129–130 ὀλολύζει ὁ γέρων (παρεπιγραφή de R)]}{\textit{Test.} Suid.~ο~194 Ὀλολύζει ὁ γέρων, εἰς τὴν θηλύτητα αὐτοῦ καὶ λαγνείαν ἀποβλέπων ; Ar.~\textit{Av.}~222 θεία μακάρων ὀλολυγή ; Aesch.~\textit{Ag.}~594–596 καὶ γυναικείῳ νόμῳ / ὀλολυγμὸν ἄλλος ἄλλοθεν κατὰ πτόλιν / ἔλασκον εὐφημοῦντες ἐν θεῶν ἕδραις ; Aesch.~\textit{Eum.}~1043 ὀλολύξατε νῦν ἐπὶ μολπαῖς ; Soph.~\textit{Trach.}~205–212 ἀνολολυξάτω δόμοις ἐφεστίοις / ἀλαλαγαῖς ἁ μελλόνυμφος, ἐν δὲ / κοινὸς ἀρσένων ἴτω / κλαγγὰ τὸν εὐφαρέτραν / Ἀπόλλω προστάταν ; Σ Ar.~\textit{Thesm.}~129 ὀλολύζει. παρεπιγραφὴ τοῦτο· εἰπόντος γὰρ τοῦ Ἀγάθωνος ὁ κηδεστὴς τοῦ Εὐριπίδου ὀλολύζει εἰς τὴν θηλύτητα αὐτοῦ καὶ μόνον οὐχὶ γυναικὸς εὐχήν ; Hsch.~ο~613 ὀλολυγή· φωνὴ γυναικῶν, ἣν ποιοῦνται ἐν τοῖς ἱεροῖς εὐχόμεναι ; Ar.~\textit{Eq.}~1327 ἀλλ᾽ ὀλολύξατε φαινομέναισιν ταῖς ἀρχαίαισιν Ἀθήναις ; Ar.~\textit{Pax}~97 μηδὲν γρύζειν ἀλλ᾽ ὀλολύζειν ; \textit{Od.}~3.450–451 αἱ δ᾽ ὀλόλυξαν / θυγατέρες τε νυοί τε καὶ αἰδοίη παράκοιτις ; \textit{Il.}~6.301 αἳ δ᾽ ὀλολυγῇ πᾶσαι Ἀθήνῃ χεῖρας ἀνέσχον}
\lp{130–132 ὡς ἡδὺ τὸ μέλος … μανδαλωτόν]}{\textit{Test.} Suid.~γ~141 Ἀριστοφάνης· ἡδὺ τὸ μέλος, ὦ πότνιαι Γενετυλλίδες, καὶ θηλυδριῶδες καὶ κατεγλωττισμένον καὶ μανδαλωτόν}
\lp{130 ὡς ἡδύ]}{Ar.~\textit{Ran.}~338 ὡς ἡδύ μοι προσέπνευσε χοιρείων κρεῶν ; Ar.~\textit{Nub.}~1025 ὡς ἡδύ σου τοῖσι λόγοις σῶφρον ἔπεστιν ἄνθος ; Ar.~\textit{Pax}~525 οἷον δὲ πνεῖς, ὡς ἡδὺ κατὰ τῆς καρδίας ; Eur.~\textit{Or.}~212 ὡς ἡδύ μοι προσῆλθες — ἐν δέοντί γε ; Eur.~fr.~443~Kannicht (\textit{Hipp. Kal.}) ὦ λαμπρὸς αἰθὴρ ἡμέρας θ᾽ ἁγνὸν φάος, / ὡς ἡδὺ λεύσσειν}
\lp{130 ὡς ἡδὺ τὸ μέλος]}{Ar.~\textit{Av.}~223–224 ὦ Ζεῦ βασιλεῦ τοῦ φθέγματος τοὐρνιθίου· / οἷον κατεμελίτωσε τὴν λόχμην ὅλην ; Ar.~\textit{Eccl.}~882–883 Μοῦσαι δεῦρ᾽ ἴτ᾽ ἐπὶ τοὐμὸν στόμα, / μελύδριον εὑροῦσαί τι τῶν Ἰωνικῶν ; Ar.~fr.~178~K.-A. (Hsch.~α~281) Ἀγαθώνειον αὔλησιν τὴν μαλακὴν … Ἀγαθὼν γὰρ ὁ τραγικὸς ἐπὶ μαλακίᾳ διεβάλλετο ; \textit{idem} Phot.~α~83 (Ἀγαθώνειος αὔλησις· ἡ μαλακή) ; Pl.~\textit{R.}~3.398e τίνες οὖν μαλακαί τε καὶ συμποτικαὶ τῶν ἁρμονιῶν; ἰαστί, ἦ δ᾽ ὅς, καὶ λυδιστὶ αὖ τινες χαλαραὶ καλοῦνται ; Plu.~\textit{Quaest.\,conv.}~3.1.1 (645e) τοῦ καλοῦ Ἀγάθωνος, ὃν πρῶτον εἰς τραγῳδίαν φασὶν ἐμβαλεῖν καὶ ὑπομῖξαι τὸ χρωματικὸν, ὅτε τοὺς Μυσοὺς ἐδίδασκεν ; D.Chr.~4.110 μύρου καὶ οἴνου ἀποπνέων ἐν κροκωτῷ … ὀρχούμενός τε καὶ ᾅδων θῆλυ καὶ ἄμουσον μέλος ; Hsch.~μ~760 μέλος· μέλημα. ἢ μέρος τοῦ σώματος. ἢ ἦχος ἡδύς, ᾆσμα, ῥυθμός ; Zen.~1.2 Ἀγαθώνιος αὔλησις: ἡ μαλακὴ, καὶ μήτε πικρὰ μήτε χαλαρὰ, ἀλλ᾽ εὔκρατος καὶ ἡδίστη. Εἴρηται δὲ ἀπὸ Ἀγάθωνος αὐλητοῦ, ὃς ἐπὶ μαλακίᾳ ἐκωμῳδεῖτο ; \textit{idem} Diogenian.~1.7 (Ἀγαθώνειος ; μαλθακὴ ; μήτε χλιαρὰ μήτε πικρὰ ; om. εὔκρατος καὶ et Εἴρηται … ἐκωμῳδεῖτο), Apostol.~1.7 (Ἀγαθώνειος ; ἡ μήτε χαλαρὰ μήτε πικρά, ἀλλ᾽ εὔκρατος καὶ ἡδίστη, μαλακὴ καὶ ἐκλελυμένη ; ἀπὸ Ἀγάθωνος αὐλητοῦ. οὗτος ἀγαθὸς ἦν τὸν τρόπον καὶ ποθεινὸς τοῖς φίλοις. ἐκωμῳδεῖτο δὲ εἰς θηλύτητα), Suid.~α~125 (μαλακὴ καὶ ἐκλελυμένη· ἢ ἡ μήτε χαλαρὰ, μήτε πικρὰ ; om. Εἴρηται … ἐκωμῳδεῖτο) ; Erasm.~\textit{Adag.}~2.4.63 Ἀγαθώνιος αὔλησις. i. e. Agathonia cantio. Recte dicetur de oratione blanda magis, quam frugifera. … Laboravit enim infamia molliciei. … Hujus multis modis Aristophanes in Thesmophoriis carminum molliciem deridet ; Sapph.~fr.~44.24–26~V. αὖλος δ᾽ ἀδυ[μ]έλης[ … ]τ᾽ ὀνεμίγνυ[το … ]ως δ᾽ ἄρα πάρ[θενοι / ἄειδον μέλος ἄγν[ον ; Pi.~\textit{O.}~7.11–12 ἄλλοτε δ᾽ ἄλλον ἐποπτεύει Χάρις ζωθάλμιος ἁδυμελεῖ / θαμὰ μὲν φόρμιγγι ; Pi.~\textit{O.}~10.93–94 τὶν δ᾽ ἁδυεπής τε λύρα / γλυκύς τ᾽ αὐλὸς ἀναπάσσει χάριν ; Hsch.~σ~1582 σταδαῖον μέλος· τὸ πρᾷον ; Hsch.~κ~3136 κνισμός· ζηλοτυπία. ἢ μέλος τι κνισμῷ αὐλούμενον ; Hsch.~κ~4236 κρουσίθυρον· μέλος τι οὕτως ἐκαλεῖτο ; Hsch.~ν~581 νηνίατος· νόμος παιδαριώδης. καὶ Φρύγιον μέλος}
\lp{130 ὦ πότνιαι Γενετυλλίδες]}{Ar.~\textit{Nub.}~51–52 ἡ δ᾽ αὖ μύρου κρόκου καταγλωττισμάτων, / δαπάνης λαφυγμοῦ Κωλιάδος Γενετυλλίδος ; Ar.~\textit{Lys.}~1–2 ἀλλ᾽ εἴ τις ἐς Βακχεῖον αὐτὰς ἐκάλεσεν, / ἢ ᾽ς Πανὸς ἢ ᾽πὶ Κωλιάδ᾽ ἢ ᾽ς Γενετυλλίδος ; [Luc.]~\textit{Am.}~42 Κωλιάδας, εἰ τύχοι, καὶ Γενετυλλίδας ἢ τὴν Φρυγίαν δαίμονα καὶ τὸν δυσέρωτα κῶμον ἐπὶ τῷ ποιμένι. τελεταὶ δὲ ἀπόρρητοι καὶ χωρὶς ἀνδρῶν ὕποπτα μυστήρια ; Paus.~1.1.5 Κωλιάδος δέ ἐστιν ἐνταῦθα Ἀφροδίτης ἄγαλμα καὶ Γενετυλλίδες ὀνομαζόμεναι θεαί· δοκῶ δὲ καὶ Φωκαεῦσι τοῖς ἐν Ἰωνίᾳ θεάς, ἃς καλοῦσι Γενναΐδας, εἶναι ταῖς ἐπὶ Κωλιάδι τὰς αὐτάς ; Σ Ar.~\textit{Thesm.}~130 δαίμων ἡ Γενετυλλὶς περὶ τὴν Ἀφροδίτην. φασὶ δὲ παρὰ τὴν γέννησιν πεποιῆσθαι τὸ ὄνομα. τινὲς δὲ περὶ τὴν Ἄρτεμίν φασιν αὐτὰς τῶν τοκετῶν ἐφόρους, καὶ πάλιν παρὰ τὴν γέννησιν τὸ ὄνομα πεποιῆσθαι ; Hsch.~γ~343 Γενετυλλίς· γυναικεία θεός, πεποιημένου τοῦ ὀνόματος παρὰ τὰς γενέσεις, ἐοικυῖα τῇ Ἀρτέμιδι. διό καὶ ταύτῃ ἐκείνας προσετίθεσαν. ἔστι δὲ ξενικὴ ἡ θεός. καὶ ἑορτή τῶν γυναικῶν ; Synag.~α~2572 Ἀφρόδιτος· ὁ Ἑρμαφρόδιτος. παραπλήσιοι δὲ τούτῳ καὶ ἄλλοι δαίμονες· Ὀρθάνης, Πρίαπος, Αἰακός, Γενετυλλίς, Τύχων, Γίγων, Κονίσαλος ; Catull.~3.1 Lugete, o Veneres Cupidinesque ; Luc.~\textit{Pseudol.}~11 εἰπὲ γάρ μοι, πρὸς πανδήμου καὶ Γενετυλλίδων καὶ Κυβήβης ; Alciphr.~2.8 ποῦ γὰρ ἐγὼ κατ᾽ ἀγρὸν ἱδρύσω Κωλιάδας ἢ Γενετυλλίδας ; Σ Ar.~\textit{Lys.}~2 ἡ Γενετυλλὶς δὲ γυναικεία θεὸς περὶ τὴν Ἀφροδίτην ἀπὸ τῆς γενέσεως [τῶν] παίδων ὠνομασμένη ; Hsch.~κ~4816 Κωλιάς· Ἀφροδίτης ἐπὶ Κωλιάδος ἐστὶν ἱερὸν ἐν τῇ Ἀττικῇ. ὁ δὲ τόπος λέγεται Κωλιάς, ἐπεὶ ἐκκείμενός ἐστιν, ὅμοιος ἀνθρώπου κώλῳ ; Hsch.~π~3156 [ποτνιάζειν· τὸ παιδικοῖς χρῆσθαι. πότνιον γὰρ τὸν δακτύλ⟨ι⟩ον λέγουσι] ; Phot.~γ~72 Γενετυλλίς· γυναικεία τις θεός. οὕτως Ἀριστοφάνης ; Suid.~κ~2240 Κωλιάς: ναός ἐστι τῆς Ἀφροδίτης οὕτω καλούμενος, ἀπὸ τοῦ συμβεβηκότος τὴν προσηγορίαν λαβών}
\lp{130 ὦ πότνιαι]}{Ar.~\textit{Thesm.}~1148–1149 ἥκετ᾽ εὔφρονες ἵλαοι, / πότνιαι, ἄλσος ἐς ὑμέτερον ; Ar.~\textit{Lys.}~742 ὦ πότνι᾽ Εἰλείθυι᾽ ἐπίσχες τοῦ τόκου ; Ar.~\textit{Lys.}~831–834 ἄνδρ᾽ ⟨ἄνδρ᾽⟩ ὁρῶ προσιόντα παραπεπληγμένον, / τοῖς τῆς Ἀφροδίτης ὀργίοις εἰλημμένον. / ὦ πότνια Κύπρου καὶ Κυθήρων καὶ Πάφου / μεδέουσ᾽, ἴθ᾽ ὀρθὴν ἥνπερ ἔρχει τὴν ὁδόν ; Ar.~\textit{Thesm.}~700 ὦ πότνιαι Μοῖραι, τί τόδε δέρκομαι ; Ar.~\textit{Eccl.}~476 καὶ ξυμφέροι γ᾽ ὦ πότνια Παλλὰς καὶ θεοί}
